\begin{center}
    \Large\textbf{\textcolor{primary}{Exponential Distribution and the Poisson Process}}
\end{center}
\vspace{0.5cm}

\section{\textcolor{secondary}{The Exponential Distribution}}
A continuous random variable $X$ is said to have an exponential distribution with parameter $\lambda$, where $\lambda > 0$, if its probability density function (PDF) is given by:
\[ f(x) = \begin{cases} \lambda e^{-\lambda x}, & x \ge 0 \\ 0, & x < 0 \end{cases} \]

Equivalently, its cumulative distribution function (CDF) is given by:
\[ F(x) = \int_{-\infty}^{x} f(y)dy = \begin{cases} 1 - e^{-\lambda x}, & x \ge 0 \\ 0, & x < 0 \end{cases} \]

\subsection{\textcolor{secondary}{Mean and Variance}}
\textbf{\textcolor{primary}{Mean ($E[X]$)}}\\
The mean is calculated using integration by parts where $u = x$, $dv = \lambda e^{-\lambda x}dx$:
\begin{align*}
E[X] &= \int_{0}^{\infty} x \lambda e^{-\lambda x} dx \\
&= \left[ -xe^{-\lambda x} \right]_{0}^{\infty} + \int_{0}^{\infty} e^{-\lambda x} dx \\
&= 0 + \left[ -\frac{1}{\lambda} e^{-\lambda x} \right]_{0}^{\infty} \\
&= \frac{1}{\lambda}
\end{align*}

\textbf{\textcolor{primary}{Variance ($Var(X)$)}}\\
First, compute $E[X^2]$ using integration by parts ($u = x^2$, $dv = \lambda e^{-\lambda x}dx$):
\begin{align*}
E[X^2] &= \int_{0}^{\infty} x^2 \lambda e^{-\lambda x} dx \\
&= \left[ -x^2 e^{-\lambda x} \right]_{0}^{\infty} + 2\int_{0}^{\infty} x e^{-\lambda x} dx \\
&= 2 \cdot \frac{1}{\lambda^2} = \frac{2}{\lambda^2}
\end{align*}

Now, compute the variance:
\begin{align*}
Var(X) &= E[X^2] - (E[X])^2 \\
&= \frac{2}{\lambda^2} - \left(\frac{1}{\lambda}\right)^2 \\
&= \frac{1}{\lambda^2}
\end{align*}

\subsection{\textcolor{secondary}{The Memoryless Property}}
A random variable $X$ is said to be without memory, or \textit{memoryless}, if:
\[ P(X > s+t \mid X > t) = P(X > s) \quad \text{for all } s, t \ge 0 \]

This is equivalent to stating:
\[ P(X > s+t) = P(X > s)P(X > t) \]

\textbf{Example:} \\
Suppose the amount of time spent in a bank is exponentially distributed with a mean of 10 minutes ($\lambda = \frac{1}{10}$). 

\textit{Probability of spending more than 15 minutes:}
\[ P(X > 15) = e^{-15\lambda} = e^{-3/2} \approx 0.223 \]

\textit{Probability of spending more than 15 minutes GIVEN she is still there after 10 minutes:}\\
Because of the memoryless property, this is simply the probability of spending at least 5 \textit{more} minutes:
\[ P(X > 5) = e^{-5\lambda} = e^{-1/2} \approx 0.607 \]

\vspace{0.5cm}

\subsection{\textcolor{primary}{Properties of Multiple Exponential Variables}}
\textbf{Proposition:} If $X_1, \dots, X_n$ are independent exponential random variables with a common rate $\lambda$, then their sum $\sum_{i=1}^{n} X_i$ is a gamma $(n, \lambda)$ random variable with density function:
\[ f(t) = \lambda e^{-\lambda t} \frac{(\lambda t)^{n-1}}{(n-1)!}, \quad t > 0 \]
\textbf{Probability of one variable being smaller than another:} \\
If $X_1$ and $X_2$ are independent with rates $\lambda_1$ and $\lambda_2$:
\begin{align*}
P(X_1 < X_2) &= \int_{0}^{\infty} P(X_1 < X_2 \mid X_1 = x)\lambda_1 e^{-\lambda_1 x} dx \\
&= \int_{0}^{\infty} e^{-\lambda_2 x} \lambda_1 e^{-\lambda_1 x} dx \\
&= \frac{\lambda_1}{\lambda_1 + \lambda_2}
\end{align*}

\section{\textcolor{primary}{The Poisson Process}}
\textbf{\textcolor{secondary}{Counting Processes:}} \\
A stochastic process $\{N(t), t \ge 0\}$ is a counting process if $N(t)$ represents the total number of ``events'' that occur by time $t$. It must satisfy:
\begin{enumerate}
    \item $N(t) \ge 0$
    \item $N(t)$ is integer-valued.
    \item If $s < t$, then $N(s) \le N(t)$.
    \item For $s < t$, $N(t) - N(s)$ equals the number of events that occur in the interval $(s, t]$.
\end{enumerate}

\textit{Independent Increments:} A counting process possesses independent increments if the numbers of events that occur in disjoint time intervals are independent. 

\vspace{0.5cm}

\textbf{\textcolor{secondary}{Little-o Notation (Definition):}} \\
The function $f(\cdot)$ is said to be $o(h)$ if:
\[ \lim_{h \to 0} \frac{f(h)}{h} = 0 \]

\textbf{\textcolor{secondary}{Definition (The Poisson Process):}} \\
The counting process $\{N(t), t \ge 0\}$ is a Poisson process with rate $\lambda > 0$ if the following axioms hold:
\begin{enumerate}
    \item $N(0) = 0$
    \item The process has independent increments.
    \item $P(N(t+h) - N(t) = 1) = \lambda h + o(h)$
    \item $P(N(t+h) - N(t) \ge 2) = o(h)$
\end{enumerate}

\vspace{0.5cm}

\subsection{\textcolor{secondary}{Interarrival Times and Sequence Distributions}}
Let $T_1$ denote the time of the first event, and $T_n$ denote the time between the $(n-1)$st and $n$th event. The sequence $\{T_n, n=1, 2, \dots\}$ is the sequence of interarrival times.

\textbf{Lemma:} \\
$T_1, T_2, \dots$ are independent and identically distributed exponential random variables with rate $\lambda$.

\textit{Proof for $T_1$:}
\[ P(T_1 > t) = P(N(t) = 0) = e^{-\lambda t} \]

Letting $S_n = \sum_{i=1}^{n} T_i$ be the time of the $n$th event, it follows from Proposition 5.1 that $S_n$ is a gamma $(n, \lambda)$ random variable:
\[ f_{S_n}(s) = \lambda e^{-\lambda s} \frac{(\lambda s)^{n-1}}{(n-1)!}, \quad s > 0 \]

\textbf{Theorem :} \\
If $\{N(t), t \ge 0\}$ is a Poisson process with rate $\lambda$, then $N(t)$ is a Poisson random variable with rate $\lambda t$. That is:
\[ P(N(t) = n) = e^{-\lambda t}\frac{(\lambda t)^n}{n!}, \quad n \ge 0 \]